\part{商预概形}\label{fga3.iii}\thispagestyle{fancy}
\begin{original}{}\label{fga3.iii-original-citation}\skipsinglepar
\par{}A. Grothendieck. "Techniques de construction et théorèmes d'existence en géométrie algébrique, III: Préschemas quotients". \emph{Séminaire Bourbaki} \textbf{13} (1960–61), 第212讲, 1961年2月. \url{http://www.numdam.org/book-part/SB_1960-1961__6__99_0/}
\end{original}
\chapter{引言}\thispagestyle{fancy}\label{fga3.iii-introduction}
\begin{remark}{}\label{fga3.iii-introduction-remark}\skipsinglepar
\par{}\emph{[Comp.]} 我们指出, \hyperref[fga3.v]{FGA~\ref*{fga3.v}} ("Picard概形: 存在性定理") 中对这里所建立理论的应用, 也可以用适当地运用Hilbert概形来替代 (参见\emph{Séminaire Mumford–Tate}, Harvard University (1961–62)). 如\cref{fga3.iii-8}所述, 这里的理论最重要的缺口是缺乏关于非固有等价关系的商的存在性判据, 例如射影群的某些作用所产生的等价关系. Mumford \cite{Mum1961} 已经在这一方向得到一个重要定理. 关于他的结果的精化及这一理论的各种应用, 见\emph{Séminaire Mumford–Tate}, Harvard University (1961–62).
\end{remark}
\par{}本讲所讨论的问题与前两讲不同: 我们试图表示的是以变动概形为自变量的某些协变函子, 而不再是反变函子. 然而, 在代数几何的许多构造问题中, 包括\hyperref[fga3.i]{[1]}和\hyperref[fga3.ii]{[2]}中的问题, 取商的过程都不可缺少. 事实上, 对于忠实平坦且拟紧的态射\(T\to  S\), 一个\(T\)-预概形\(X\)上的\emph{下降数据的有效性}问题, 等价于\(X\)关于该下降数据在\(X\)上所定义的平坦等价关系的商的存在性问题, 其中要求这个商具有下面将研究的一些合理性质; \hyperref[fga3.i-a.2.c]{[1], A, §2.c}中提出的问题, 或许可以与本讲\cref{fga3.iii-2}所提出的问题一并解决. 同样, 一个\(S\)-概形\(X\)的\emph{Picard概形} (定义见\hyperref[fga3.ii-c.3]{[2], C, §3}) 可以用多种方式定义为某些其他概形 (正除子概形, 或射影嵌入概形) 关于平坦等价关系的商, 而这些辅助概形的定义与构造也较为简单: 它们基本上是\(\operatorname {Hom}_S(X,Y)\)型的概形以及\hyperref[fga3.ii-c.2]{[2], C, §2}中定义的变体; 下一讲将讨论它们的构造, 其中须加上适当的射影性假设. 因而, 将本讲的结果与下一讲的结果结合, 就能在适当假设下构造Picard概形.
\par{}预概形中的取商问题仍有尚未解决的疑问. 最重要的一个将在\cref{fga3.iii-8}中说明. 目前, 它仍是构造\emph{任意亏格曲线在整数环上的模概形}、\emph{极化阿贝尔簇的模概形}等的唯一障碍. 因此, 它的解决值得代数群专家为之努力.
\chapter{等价关系与有效等价关系}\thispagestyle{fancy}\label{fga3.iii-1}
\par{}设\(\mathcal {C}\)是一个范畴, \(X\)是\(\mathcal {C}\)的一个对象. 如果对于\(\mathcal {C}\)的每个对象\(T\), 相应的映射\[     p_1(T),p_2(T)\colon  R(T)\rightrightarrows  X(T)   \] (这里对\(\mathcal {C}\)的任意对象\(Y\), 置\(Y(T)=\operatorname {Hom}(T,Y)\)) 所定义的映射\[     R(T)\to  X(T)\times  X(T),\] 都给出\(R(T)\)到集合\(X(T)\)上某个等价关系的图的双射, 则称态射对\[     p_1,p_2\colon  R\rightrightarrows  X,   \]是\(\mathcal {C}\)中以\(X\)为\emph{目标}、以\(R\)为\emph{来源}的"\emph{等价对}". 在以\(X\)为目标的等价对之间有一个显然的等价关系; 它的一个等价类称为\(X\)上的一个\emph{\(\mathcal {C}\)-等价关系}, 或在不会混淆时简称等价关系.
\par{}如果\(X\times  X\)存在, 则给出\(X\)上的一个等价关系, 等价于给出\(X\times  X\)的一个子对象\(R\), 使得对于\(\mathcal {C}\)的每个对象\(T\), 与\(R(T)\)相应的\((X\times  X)(T)=X(T)\times  X(T)\)的子集是\(X(T)\)上一个等价关系的图. 将投影\(\mathrm {pr}_1\)、\(\mathrm {pr}_2\)分别诱导的从\(R\)到\(X\)的态射记作\(p_1\)、\(p_2\), 上述条件就是说\((p_1,p_2)\)是一个等价对. 按照\hyperref[fga3.ii-a.1]{[2], A, §1}中阐述的一般原则, 也可以在\(\mathcal {C}\)中以交换图表示\(X(T)\)中的\(R(T)\)所满足的集合论等价关系公理 (这里假设\(X\times  X\)和纤维积\((R,p_2)\times _X(R,p_1)\)均存在). 我们不需要这样做.
\par{}只要有同来源\(R\)、同目标\(X\)的一对态射\((p_1,p_2)\), 就可以将这一对的\emph{余核}定义为\(\mathcal {C}\)的一个对象\(Y\), 它表示关于\(Z\)的协变函子\[     \operatorname {Hom}_{p_1,p_2}(X,Z);\] 这里的值是所有满足\(up_1=up_2\)的从\(X\)到\(Z\)的态射\(u\)所组成的集合. 如果\(Y\)存在, 则它在相差唯一同构的意义下确定. 我们将它记作\(Y/(p_1,p_2)\), 或滥用记号记作\(Y/R\); 后一个记号主要在\((p_1,p_2)\)为等价对时使用, 此时习惯于在记号上把该对定义的等价关系与\(R\)混同. 注意, 若将\(Y\)视为\(X\)的商, 那么它确实只依赖于态射对\((p_1,p_2)\)定义的等价关系.
\par{}现在从一个态射\[     f\colon  X\to  Y   \]出发, 它使我们能把\(X\)看成"\(Y\)之上的对象", 并假设纤维积\[     \mathcal {R}(f)     = X\times _Y X   \]存在. 设\(p_1\)、\(p_2\)是它的两个投影. 则\((p_1,p_2)\)是一个等价对, 称为与态射\(f\)\emph{相伴}的等价对. 它因而定义了一个等价关系, 称为与态射\(f\)\emph{相伴}的等价关系.
\par{}称以\(X\)为目标、\(R\)为来源的态射对\((p_1,p_2)\)为\emph{有效等价对}, 如果满足下列条件: 
\begin{enumerate}
    
 \item[i.]余核\(Y=X/(p_1,p_2)\)存在;

    
 \item[ii.]纤维积\(X\times _Y X\)存在;

    
 \item[iii.]以\(p_1\)和\(p_2\)为分量的态射\(R\to  X\times _Y X\)为同构.

  
\end{enumerate} 此时\((p_1,p_2)\)确实是等价对. 我们也称它定义的等价关系为\emph{有效等价关系}.
\par{}称态射\(f\colon  X\to  Y\)为\emph{有效满态射}, 如果满足下列条件: 
\begin{enumerate}
    
 \item[i.]纤维积\(R=X\times _Y X\)存在;

    
 \item[ii.]商\(X/(p_1,p_2)\)存在, 这里\(p_1\)、\(p_2\)是从\(R\)到\(X\)的两个投影;

    
 \item[iii.]由\(f\)诱导的态射\(X/(p_1,p_2)\to  Y\)为同构.

  
\end{enumerate} 此时\(f\)确实是满态射, 甚至是严格满态射 (参见\hyperref[fga3.i-a.2.c]{[1], A, §2.c}); 如果纤维积\(X\times _Y X\)存在, 则反之亦然. 我们也称满态射\(f\)所定义的\(X\)的商对象为\(X\)的一个\emph{有效商}.
\par{}上述定义蕴涵下述"\emph{Galois对应}":
\begin{proposition}{1.1}\setnumber{1.1}\label{fga3.iii-1-proposition-1.1}\skipsinglepar
\par{}\(X\)上有效等价关系\(R\)的集合与\(X\)的有效商\(Y\)的集合之间, 存在一个保持自然序的双射对应: \(R\)对应于有效商\(X/R\), 而\(Y\)对应于典范投影\(X\to  Y\)所定义的有效等价关系, 即由带有两个投影的纤维积\(X\times _Y X\)所定义的关系.
\end{proposition}
\par{}在性质很好的范畴 (集合、集合层等) 中, 每个商都是有效的, 每个等价关系也都是有效的. 对于给定预概形\(S\)之上的预概形范畴等范畴, 这就不再成立, 即使\(S\)是一个域的谱, 甚至只考虑\(S\)上的有限概形, 也不成立. 有效性问题, 乃至商的存在性问题 (在\(S\)上的预概形不有限时), 往往都很微妙.
\chapter{例: \(S\)上的有限预概形}\thispagestyle{fancy}\label{fga3.iii-2}
\par{}设\(\mathcal {C}\)是\(S\)上的有限预概形范畴, 其中假设基底局部Noether. 那么\(\mathcal {C}\)等价于\(S\)上凝聚交换代数层范畴的反范畴; 如果\(S\)是环为\(A\)的仿射概形, 则它等价于在\(A\)上有限的\(A\)-代数 (即在\(A\)上作为模为有限型的代数) 范畴的反范畴. 因而立即得到, \(\mathcal {C}\)中的有限投射极限与有限归纳极限均存在. 对前者, 这是熟知的, 而且不需任何有限性假设: 预概形\(X\)和\(Y\)在\(S\)上的纤维积对应于相应代数的张量积\(B\otimes _A C\), 由两个\(A\)-代数同态\(u,v\colon  C\rightrightarrows  B\)定义的两个态射\(X\rightrightarrows  Y\)的核, 对应于\(B\)模掉由\(u(c)-v(c)\)生成的理想所得的商, 等等. 对于有限归纳极限, 只需一方面考虑有限和, 它们对应于\(A\)-代数的通常直积; 另一方面考虑态射对\(X\rightrightarrows  Y\)的余核, 它对应于\(C\)中使两个同态\(u,v\colon  C\rightrightarrows  B\)取值相同的元素组成的子环, 这一点可以直接看出 (由Noether假设, 此子环在\(A\)上有限). 还应指出, 利用Noether假设, 可以证明在\(S\)上的有限预概形范畴\(\mathcal {C}\)中如此构造的有限归纳极限, 特别是商, 实际上也是\emph{所有}预概形范畴中的商.
\par{}如\hyperref[fga3.i]{[1]}所述, \emph{\(\mathcal {C}\)中存在非有效的满态射} (也即非严格的满态射; 因为纤维积存在, 这两者相同). 如果不加平坦性假设, \emph{我不知道等价关系是否总是有效的}. 在这一方向, 我只得到一些很局部的肯定结果, 而这些结果对于形式模理论基本定理的证明不可缺少 (参见\hyperref[fga3.ii-b-theorem-1]{[2], B, 定理1}). 注意, 在所提问题中, 很容易归结到\(S\)是剩余域代数闭的Artin局部环的谱的情形. 但即使\(A\)是一个域, 答案也还不知道.
\par{}也可以考虑\(S\)上的预概形\(X\)不再假设在\(S\)上有限, 而\(X\)上的一个等价关系\(R\)满足\(p_1\colon  R\to  X\)是有限态射的情形. 此时称\(R\)为\emph{有限等价关系}. 为简单起见, 假设\(S\)和\(X\)仿射 (于是\(R\)也仿射, 从而问题归结为纯交换代数的情形), \emph{即使在这种情形, 我们也不知道商\(X/R=Y\)是否存在, 以及典范态射\(X\to  Y\)是否有限}. (最简单的情形是假设\(S\)是域\(k\)的谱, 而\(X\)是\(k[t]\)的谱, 即仿射直线.) 当然, 如果上述两个问题的答案都为肯定, 那么在所述情形就能断定\(R\)是有效的. 注意, 若不是\(X\)中的等价关系图, 而只有\cref{fga3.iii-4}意义下\(X\)中的一个预等价关系图, 则商\(Y\)的\emph{存在性}问题和\(f\colon  X\to  Y\)的\emph{有限性}问题仍可用完全相同的方式提出.
\par{}沿某些闭子预概形"粘合"给定的预概形\(X_i\), 以构造预概形时, 会遇到关于一般有限等价关系的取商问题; 粘合法则恰好由\(X_i\)之和构成的预概形\(X\)上的有限等价关系表示. 也可以预期, 这些问题及其许多变体的解决, 将是沿\hyperref[fga3.ii]{[2]}开创的方向建立一般非射影构造方法的先决条件.
\par{}作者所知道的唯一一般性的肯定事实如下:
\begin{proposition}{2.1}\setnumber{2.1}\label{fga3.iii-2-proposition-2.1}\skipsinglepar
\par{}设\(S\)是局部Noether预概形, \(s\)是\(S\)的一点, \(\Omega \)是\(k(s)\)的一个代数闭扩张. 考虑相应的"纤维函子"\(F\), 它把每个在\(S\)上有限的\(S\)-概形\(X\)对应于\(X/S\)的所有\(\Omega \)值点组成的集合. 此函子显然左正合, 而且是\emph{右正合}的, 即它与有限归纳极限交换, 特别是与态射对的余核交换.
\end{proposition}
\par{}对\(S\)的所有"几何点"应用这个结果, 就得到\(\mathcal {C}\)的"商"范畴\(\mathcal {C}'\) (即"模掉满的根态射"来推理, 也就是形式地给这类态射添上逆) 是一个"几何"范畴, 即满足与集合范畴相同的"有限性质". 特别是每个等价关系都有效. 由此得出: 若\(R\)是\(X\)上的等价关系, 且\(X\)在\(S\)上有限, 则典范态射\(R\to  X\times _Y X\) (其中\(Y=X/R\)) 是\emph{根且满}的; 事实上, 因为它是单态射, 所以是满的闭浸入.
\chapter{群作用的情形}\thispagestyle{fancy}\label{fga3.iii-3}
\par{}现在重新假设\(\mathcal {C}\)是任意范畴. 设\(G\)、\(X\)是\(\mathcal {C}\)的对象, 且\(G\)是作用于对象\(X\)的\(\mathcal {C}\)-群. 这就是说 (参见\hyperref[fga3.ii-a.1]{[2], A, §1}), 对\(\mathcal {C}\)的每个对象\(T\), 在\(G(T)\)上给定一个群结构, 并在\(X(T)\)上给定以\(G(T)\)为算子群的集合结构, 使得当\(T\)变动时, 这些结构"随\(T\)函子地变动". 若\(G\times  G\)和\(G\times  X\)在\(\mathcal {C}\)中存在, 则这种结构也可以定义为一对态射\[     
\begin {gathered}       G\times  G\to  G     \\\pi \colon  G\times  X\to  X     
\end {gathered},\] 满足对于\(\mathcal {C}\)的每个对象\(T\), 集合\(G(T)\)和\(X(T)\)上的相应合成律使\(G(T)\)成为作用于\(X(T)\)的群. 将这个公理翻译为\(\mathcal {C}\)中若干图的交换性是容易的, 但很繁琐, 而且在我所知的所有情形中实际上完全没有必要.
\par{}假设\(G\times  X\)存在, 考虑两个态射\[     p_1,p_2\colon  G\times  X\rightrightarrows  X,\] 其中\[     
\begin {aligned}       p_1 &= \mathrm {pr}_1     \\p_2 &= \pi .     
\end {aligned}.\] 立即可见, \((p_1,p_2)\)是等价对当且仅当对于\(\mathcal {C}\)的每个对象\(T\), 该对定义的映射\[     G(T)\times  X(T)     \sim  (G\times  X)(T) \to  X(T)\times  X(T)   \]是单射, 即群\(G(T)\)在集合\(X(T)\)上\emph{自由地}作用; 换言之, 若\(g\in  G(T)\)、\(x\in  X(T)\)且\(g\cdot  x=x\), 则\(g\)是群\(G(T)\)的单位元. 此时称\(G\)在\(X\)上\emph{自由作用} (或称\(X\)是\emph{\(G\)下的主\(\mathcal {C}\)-空间}). 于是与\((p_1,p_2)\)相伴的等价关系称为在\(X\)上自由作用的\emph{群\(G\)所定义的等价关系}. 如果\(X\times  X\)也存在, 并考虑由\((p_1,p_2)\)定义的态射\[     p\colon  G\times  X\to  X\times  X,\] 那么\(G\)自由作用的条件就是说\(p\)是一个\emph{单态射}.
\par{}当然, 即使\(G\)在\(X\)上的作用不是自由的, 我们仍然希望有\(X\)关于\(G\)的商的存在性判据, 即上述态射对\((p_1,p_2)\)的余核的存在性判据.
\par{}这个余核通常记作\(X/G\); 若\(G\)在左边作用, 则更适合记作\(G\backslash  X\) (前一个记号留给\(G\)在右边作用的情形). 应当指出, 即使\(G\times  X\)在\(p\)下的"像"存在, 记作\(R\) (例如将此像定义为\(X\times  X\)中能使\(p\)经由它分解的最小子对象), 它通常也不是\(X\)上的等价关系. 若此时试图直接关于\(R\)取商 (更准确地说, 关于由两个投影\(\mathrm {pr}_i\)诱导的从\(R\)到\(X\)的态射对取商), 就会丢失原态射对\((p_1,p_2)\)的特有性质. 因此有必要推广等价关系的概念, 使它能直接适用于由一个带作用的\(\mathcal {C}\)-群定义的态射对.
\chapter{预等价关系}\thispagestyle{fancy}\label{fga3.iii-4}
\par{}回顾一下, \emph{群胚}是指所有态射均为同构的范畴. 另一方面, 一个范畴由两个基本集合\(X\)、\(R\) (前者是\emph{对象}集, 后者是\emph{箭头}集) 以及下列结构组成: 
\begin{enumerate}
    
 \item[i.]一对映射\[         p_1,p_2\colon  R\rightrightarrows  X,\] 称为\emph{来源映射}和\emph{目标映射};

    
 \item[ii.]一个映射\[         \pi \colon (R,p_2)\times _X(R,p_1) \to  R,\] 称为\emph{合成映射}.

  
\end{enumerate} 这些数据应满足熟知的公理, 这里不再重复. 公理可以表示为某些图的交换性, 再加上一个使另外两个图交换的必然唯一的映射\(D\colon  X\to  R\)的存在性; \(D\)把对象对应于其恒等映射, 并满足\[     p_1\circ  D = p_2\circ  D = \operatorname {id}_X.\] 说这个范畴是群胚, 就意味着存在一个必然唯一的映射\[     s\colon  R\to  R,\] 称为\(R\)的\emph{对称}, 它把每个箭头变成其逆箭头. 这可以表示为由\(s\)、\(\Delta \)及上述数据构成的另外四个图的交换性, 其中前两个可写成\[     
\begin {aligned}       p_1\circ  s &= p_2     \\p_2\circ  s &= p_1.     
\end {aligned}.\]
\par{}回顾这些概念之后, \hyperref[fga3.ii-a.1]{[2], A, §1}中的一般定义特别说明了, 对于任意范畴\(\mathcal {C}\)的一对对象\((X,R)\), 所谓"\emph{\(\mathcal {C}\)-范畴}结构" (分别为\emph{\(\mathcal {C}\)-群胚}结构) 应是什么意思: 按定义, 就是对\(\mathcal {C}\)的每个对象\(T\), 在以\(X(T)\)为对象集、以\(R(T)\)为箭头集的集合对上给定一个集合论意义的范畴 (分别为群胚) 结构, 并使这些结构"随\(T\)函子地变动". 于是得到两个态射\[     p_1,p_2\colon  R\rightrightarrows  X,\] 称为\emph{来源态射}和\emph{目标态射}, 并且若有关的纤维积存在, 则还有一个态射\[     \pi \colon  (R,p_2)\times _X(R,p_1) \to  R,\] 称为\emph{合成态射}. 这三个态射就足以确定\((X,R)\)上的范畴 (分别为群胚) 结构, 对它们所加的条件是: 对于每个\(T\), 与\(X(T)\)和\(R(T)\)相应的三个映射在集合对\((X(T),R(T))\)上定义一个范畴 (分别为群胚) 结构. 必要时, 也可将此表示为某些图的交换性, 其中涉及一个完全确定的态射\[     D\colon  X\to  R,\] 以及在群胚情形中的一个完全确定的态射\[     s\colon  R\to  R;\] 这些图与"集合论"情形相同. 幸而实际上不需要把公理作如此繁琐的解释. 能够用某些纤维积之间的态射及态射等式来表示这些数据和公理, 在理论上的唯一意义是: 若有一个左正合函子\(F\colon \mathcal {C}\to \mathcal {C}'\) (即与有限直积和纤维积交换的函子), 则它把\(\mathcal {C}\)-范畴 (分别为\(\mathcal {C}\)-群胚) 变成\(\mathcal {C}'\)-范畴 (分别为\(\mathcal {C}'\)-群胚), 条件是有限直积和纤维积在\(\mathcal {C}\)中存在.
\par{}实际中重要的是, 应懂得如何把态射\(p_1\)、\(p_2\)、\(\pi \)、\(D\)和\(s\)理解为\(\mathcal {C}\)中一个适当的半单纯或单纯对象的\emph{单纯运算} (至少在\(\mathcal {C}\)中纤维积存在时如此). 为固定术语, 引入\emph{单纯形型}范畴\(\mathcal {S}\): 它的对象为形如\[     \Delta _n     = [0,n]   \]的有限集合, 其中\(n\in \mathbb {Z}_{\geqslant 0}\), 而\([0,n]\)表示从\(0\)到\(n\)的整数闭区间; 它的态射为这些有限集合之间的\emph{任意映射}. 注意\(\mathcal {S}\)等价于\emph{非空}有限集合及其间所有映射构成的范畴. \(\mathcal {S}\)中显然存在\emph{非空}有限族对象的和, 也存在两个对象沿第三个对象的粘合和 (它是纤维积的对偶运算). 以\(\mathcal {S}'\)表示\(\mathcal {S}\)的具有相同对象的子范畴, 但态射只取\(\Delta _n\)之间的\emph{递增映射}. 此范畴等价于非空有限全序集的范畴. 在这个范畴中, 两个对象的和从不存在, 而两个对象\(A\)、\(B\)沿第三个对象\(C\)的粘合和一般也不存在 (例如取\(C=\Delta _0\)、\(A=B=\Delta _1\), 且两个结构映射\(u\colon  C\to  A\)与\(v\colon  C\to  B\)相同). 不过, 在某些情形粘合和\emph{确实}存在; 考虑\[     
\begin {gathered}       A = \Delta _m       \qquad  B = \Delta _n       \qquad  C = \Delta _0     \\u(0) = m       \qquad  v(0) = 0     
\end {gathered},\] 此时有\[     A\coprod _C B     = \Delta _{m+n}.\]
\par{}范畴\(\mathcal {C}\)中的\emph{单纯对象} (分别为\emph{半单纯对象}) 定义为从\(\mathcal {S}\) (分别为\(\mathcal {S}'\)) 到\(\mathcal {C}\)的一个反变函子\(K\). 因此单纯对象通过限制定义一个半单纯对象, 但前者与后者的本质区别在于\(K_n=K(\Delta _n)\)中还有\emph{对称运算}, 它们对应于\(n+1\)个元素的对称群的元素在函子\(K\)下的像 (此群视为\(\mathcal {S}\)中\(\Delta _n\)的自同构群).
\par{}按上述约定, 对每个\(n\), 以\(\Delta '_n\) (分别为\(\Delta ''_n\)) 表示对象集为\(\Delta _n\)、箭头集由\(\Delta _n\)上的"混沌序"关系 (分别为自然全序关系) 定义的有限范畴 (即箭头集是这一关系的图). 显然, \(\Delta '_n\) (分别为\(\Delta ''_n\)) 函子地依赖于\(\mathcal {S}\) (分别为\(\mathcal {S}'\)) 的对象\(\Delta _n\). 因而若\(Z\)是一个范畴, 则当\(\Delta _n\)变动时, \(\operatorname {Hom}(\Delta '_n,Z)\) (分别为\(\operatorname {Hom}(\Delta ''_n,Z)\)) 是从范畴\(\mathcal {S}\) (分别为\(\mathcal {S}'\)) 到集合范畴的一个函子, 即所谓与\emph{范畴\(Z\)相伴}的\emph{单纯集} (分别为\emph{半单纯集}), 记作\(Z'\) (分别为\(Z''\)). 此外, 从与\(Z'\)相应的半单纯集到\(Z''\)有一个显然的自然同态, 它是同构当且仅当\(Z\)是群胚. 于是:
\begin{proposition}{4.1}\setnumber{4.1}\label{fga3.iii-4-proposition-4.1}\skipsinglepar
\par{}从范畴的范畴到半单纯集范畴的函子\(Z\mapsto  Z''\)完全忠实, 并在\emph{范畴}的范畴与满足以下条件的半单纯集范畴之间定义等价: 即从\(\mathcal {S}'\)到\(\mathtt {Set}\)的反变函子\(K\)\emph{把上述类型的粘合和\(A\coprod _C B\)变成集合的纤维积}.
\par{}同样, 从群胚范畴到单纯集范畴的函子\(Z\mapsto  Z'\)完全忠实, 并在\emph{群胚}的范畴与满足以下条件的单纯集范畴之间定义等价: 即从\(\mathcal {S}\)到\(\mathtt {Set}\)的反变函子\(K\)\emph{把粘合和变成纤维积}.
\end{proposition}
\par{}因此可以把范畴看成半单纯集的特殊情形, 把群胚看成单纯集的特殊情形; 当然须在"相差同构"的意义下这样说, 这正是在用一种结构解释另一种结构时的严格说法. 通常的归结到集合论情形的办法于是给出:
\begin{corollary}{4.2}\setnumber{4.2}\label{fga3.iii-4-corollary-4.2}\skipsinglepar
\par{}将范畴、群胚和单纯集分别换成\(\mathcal {C}\)-范畴、\(\mathcal {C}\)-群胚和\(\mathcal {C}\)-单纯对象, 上述断言仍然成立, \emph{条件是}\(\mathcal {C}\)中存在纤维积.
\end{corollary}
\par{}与\(\mathcal {C}\)中的范畴\((X,R,\ldots )\)相伴的\(\mathcal {C}\)中的半单纯对象\(K\), 可以具体写出: 将\(K\)的分量\(K_n=K(\Delta _n)\)视为\((R,p_1)\)在\(X\)上的\((n+1)\)重纤维积, 或者更方便地用递推公式\[     
\begin {aligned}       K_0 &= R     \\K_n &= (K_{n-1},p_n^{(n-1)})\times _X(R,p_1)     
\end {aligned}   \]来定义; 这里\(p_i^{(n-1)}\) (其中\(0<i<n-1\)) 是从\(K_{n-1}\)到\(X\)的自然投影 (也可递推定义). 如此, \(p_1\)、\(p_2\)、\(\pi \)、\(D\)和\(s\)便可分别理解为与\(\mathcal {S}\)中下列态射相应的单纯运算: \(\Delta _1\)的\(0\)面、\(\Delta _1\)的\(1\)面、\(\Delta _2\)的\((0,2)\)面、退化\(\Delta _1\to \Delta _0\)以及\(\Delta _1\)的对称. 其余所有半单纯 (分别为单纯) 运算均可由上述四个 (分别为五个) 运算通过合成和纤维积形式地得到.
\par{}现在将一个范畴的对象\(X\)上的\emph{预等价关系}定义为以\(X\)为对象之对象的一个群胚的数据. 这样的数据, 除其他部分以外, 还给出对象\(R\)及两个态射\[     p_1,p_2\colon  R\rightrightarrows  X.\] 但应注意, 与等价对的情形相反, 仅这些数据并不能确定所讨论的结构. 本讲之所以研究这一概念, 是为了得到能够取商的判据, 即能够作出态射对\((p_1,p_2)\)的余核的判据. 因而这个问题的陈述不涉及群胚固有的其他数据. 不过, 在下面各结果的证明中, 我们将用到这些额外数据, 特别是直到\(3\)维的单纯运算 (包括对称运算), 其中将出现\(R\)在\(X\)上的四重纤维积.
\par{}\(\mathcal {C}\)的对象\(X\)上的一个等价关系定义一个预等价关系: 只需在集合论情形说明这一点, 此时对集合\(X\)上的一个等价关系, 配以对象集为\(X\)、箭头集为该等价关系的图的集合的群胚即可.
\par{}作用于\(\mathcal {C}\)的对象\(X\)的\(\mathcal {C}\)-幺半群\(G\), 定义一个以\(R=G\times  X\)和\(X\)为基本对象的\(\mathcal {C}\)-范畴 (条件是\(G\times  X\)存在); 而该范畴是\(\mathcal {C}\)-群胚当且仅当\(G\)是群. 仍只须在集合论情形证明. 此时把箭头\((g,a)\)与\((g',g\cdot  a)\)的合成定义为\[     (g',g\cdot  a) \circ  (g,a)     = (g'g,a);\] 也就是说, 若\(a,b\in  X\), 则按定义, \(\operatorname {Hom}(a,b)\)是把\(a\)变到\(b\)的搬运子, 而态射依靠\(G\)中元素的合成而合成.
\begin{remark}{}\label{fga3.iii-4-remark}\skipsinglepar
\par{}只要暗中假设所有涉及的对象均属于某个固定的"宇宙" (它本身是一个集合), 就可避开\hyperref[fga3.iii-4-proposition-4.1]{命题~\ref*{fga3.iii-4-proposition-4.1}}这类陈述所引起的逻辑困难.
\end{remark}
\chapter{关于有限平坦等价关系的商}\thispagestyle{fancy}\label{fga3.iii-5}
\begin{theorem}{5.1}\setnumber{5.1}\label{fga3.iii-5-theorem-5.1}\skipsinglepar
\par{}设\(X=\operatorname {Spec}(B)\)是仿射概形, \(\mathcal {R}\)是\(X\)上的一个预等价关系, 且其分量\(R_1\)仿射, 记为\(R_1=\operatorname {Spec}(C)\). 假设第一投影\(p_1\colon  R_1\to  X\)是有限局部自由态射, 即相应的环同态\(p'_1\colon  B\to  C\)使\(C\)成为有限型投射\(B\)-模. 设\(A\)是\(B\)中由同态对\(p'_1,p'_2\colon  B\rightrightarrows  C\)的核给出的子环 (即满足\(p'_1(b)=p'_2(b)\)的元素\(b\)的集合). 置\(Y=\operatorname {Spec}(A)\), 以\(f\colon  X\to  Y\)表示由\(A\)嵌入\(B\)定义的态射. 此时有: 
\begin{enumerate}
      
 \item[i.]\(B\)在\(A\)上整, 即\(f\)是整态射.

      
 \item[ii.]态射\(f\)是满的, 其纤维是\(X\)模\(\mathcal {R}\)的集合论等价类\(p_2(p_1^{-1}(x))\), 而\(Y\)的拓扑是\(X\)的拓扑的商拓扑.

      
 \item[iii.]\(Y\)是预概形范畴中\(X\)关于\(\mathcal {R}\)的商.

      
 \item[iv.]若\(\mathcal {R}\)来自一个等价\emph{关系}, 则态射\(f\colon  X\to  Y\)有限局部自由 (即\(B\)是有限型投射\(A\)-模), 且该等价关系有效, 即\(R_1\to  X\times _Y X\)是同构.

    
\end{enumerate}
\end{theorem}
\par{}这个定理推广了有限群\(G\)以自同构作用于环\(B\)、\(A\)为不变元环时的熟知定理, 其证明与已知的证明类似. 还可如下精化(iii):
\begin{corollary}{5.2}\setnumber{5.2}\label{fga3.iii-5-corollary-5.2}\skipsinglepar
\par{}典范态射\(R_1\to  X\times _Y X\)是\emph{满}的.
\end{corollary}
\par{}仍设\(\mathcal {R}\)是\(X\)上一个"有限局部自由"的预等价关系, 但现在\(X\)是任意预概形. 假设能找到预概形\(Y\)及态射\(f\colon  X\to  Y\), 使\(fp_1=fp_2\), 并使\(Y\)上的环层同态序列\[     \mathcal {O}_Y \to  f_*(\mathcal {O}_X) \rightrightarrows  g_*(\mathcal {O}_R)   \]正合 (其中\(g=fp_i\)). 由上述定理可得与该定理(i)至(iv)相应的结论; 特别由(iii), \(Y\)是\(X\)关于\(\mathcal {R}\)的商, 因而在相差唯一同构的意义下确定. 在这些条件下, 称\(X\)上的预等价关系\(\mathcal {R}\)为\emph{容许}的. 按此定义:
\begin{theorem}{5.3}\setnumber{5.3}\label{fga3.iii-5-theorem-5.3}\skipsinglepar
\par{}设\(X\)是预概形, \(\mathcal {R}\)是\(X\)上的预等价关系, 且\(p_1\colon  R_1\to  X\)有限局部自由. \(\mathcal {R}\)容许的充要条件是\(X\)模\(\mathcal {R}\)的每个集合论等价类\(p_2(p_1^{-1}(x))\)都包含在一个仿射开子集中 (如果\(X\)的每个有限子集都包含于一个仿射开子集, 这个条件总能满足, 例如\(X\)在某个仿射概形上拟射影时).
\end{theorem}
\par{}事实上容易证明, \(X\)模\(\mathcal {R}\)的每个等价类此时均包含于一个在\(\mathcal {R}\)下\emph{稳定}的仿射开子集, 于是将应用\hyperref[fga3.iii-5-theorem-5.1]{定理~\ref*{fga3.iii-5-theorem-5.1}}得到的各块粘合起来, 就构造出商\(Y\).
\begin{corollary}{5.4}\setnumber{5.4}\label{fga3.iii-5-corollary-5.4}\skipsinglepar
\par{}假设这个条件满足, 并且\(\mathcal {R}\)来自一个等价关系. 则该等价关系有效, 即\(R_1\to  X\times _Y X\)是同构, 且\(f\colon  X\to  Y\)是有限局部自由态射.
\end{corollary}
\par{}于是由下降立即得到:
\begin{corollary}{5.5}\setnumber{5.5}\label{fga3.iii-5-corollary-5.5}\skipsinglepar
\par{}在\hyperref[fga3.iii-5-corollary-5.4]{推论~\ref*{fga3.iii-5-corollary-5.4}}的条件下, \(X\)在\(Y\)上处处秩为\(n\)的充要条件是\((R_1,p_1)\)在\(X\)上处处秩为\(n\). 若\(X\)和\(R_1\)是\(Z\)-预概形, 且\(p_1\)、\(p_2\)是\(Z\)-态射 (因而\(Y\)是\(Z\)-预概形), 则\(X\)在\(Z\)上平坦当且仅当\(Y\)在\(Z\)上平坦.
\end{corollary}
\par{}归纳如下:
\begin{scholium}{}\label{fga3.iii-5-scholium}\skipsinglepar
\par{}给出预概形间一个有限、局部自由且满的态射\(f\colon  X\to  Y\), 等价于给出一个预概形\(X\)及其上的等价关系\(R\), 满足\(p_1\colon  R\to  X\)有限局部自由, 且每个等价类\(p_2(p_1^{-1}(x))\)包含于一个仿射开子集.
\end{scholium}
\begin{remarks}{5.6}\setnumber{5.6}\label{fga3.iii-5-remarks-5.6}\skipsinglepar
\begin{enumerate}\item{}我们不需要作任何Noether性假设.

    \item{}这种取商方式包含Cartier的"不可分下降"作为特例: 它相应于确定有限局部自由态射\(f\colon  X\to  Y\), 使得\(f_*(\mathcal {O}_X)\)相对于\(\mathcal {O}_Y\)具有\(p\)-基 (这里\(X\)是给定预概形, 其层\(\mathcal {O}_X\)被素数\(p>0\)零化). 注意, 不对局部环作任何正则性假设, 也不假设\(X\)是域上的代数概形, 也很容易表述这一结果. 取\(X\)为特征\(p\)的一个域的谱, 就得到Jacobson–Bourbaki理论.

    \item{}Gabriel在域\(k\)上有限交换群的取商理论中已经得到\hyperref[fga3.iii-5-theorem-5.3]{定理~\ref*{fga3.iii-5-theorem-5.3}}的一个特例. (比较\hyperref[fga3.iii-7-corollary-7.4]{推论~\ref*{fga3.iii-7-corollary-7.4}}.)
\end{enumerate}
\end{remarks}
\chapter{关于固有平坦等价关系的商}\thispagestyle{fancy}\label{fga3.iii-6}
\begin{theorem}{6.1}\setnumber{6.1}\label{fga3.iii-6-theorem-6.1}\skipsinglepar
\par{}设\(S\)是\emph{局部Noether}预概形, \(X\)是\emph{拟射影}的\(S\)-概形, \(\mathcal {R}\)是\(X\)上的预等价关系, 满足: 
\begin{enumerate}
      
 \item[a.]\(p_1\colon  R_1\to  X\)固有且平坦;

      
 \item[b.]\(R_1\to  X\times _S X\)是有限态射 (由(a), 等价地说, 是纤维有限的态射; 若\(\mathcal {R}\)来自等价关系, 则此条件自动满足).

    
\end{enumerate} 在这些条件下: 
\begin{enumerate}
      
 \item[i.]\(Y=X/\mathcal {R}\)存在, 且当\(S\)为Noether时, 它在\(S\)上拟射影. (\emph{[Trans.] \emph{[Comp.]} 关于\(Y=X/\mathcal {R}\)在\(S\)上拟射影的事实, 目前只在\(\mathcal {R}\)来自等价关系的情形得到证明.})

      
 \item[ii.]典范态射\(f\colon  X\to  Y\)满、固有且开, 其纤维是\(X\)模\(\mathcal {R}\)的等价类\(p_2(p_1^{-1}(x))\), 因而\(Y\)可与\(X\)关于\(\mathcal {R}\)所定义的集合论等价关系的拓扑商空间认同. 最后, \(R_1\to  X\times _Y X\)是满的.

      
 \item[iii.]若\(\mathcal {R}\)来自等价关系, 则该等价关系有效, 即\(R_1\to  X\times _Y X\)是同构; 而且\(f\colon  X\to  Y\)平坦, 从而忠实平坦.

    
\end{enumerate}
\begin{proof}
\par{}证明时, 可考虑\(X\)关于\(\mathcal {R}\)的适当拟截面, 将问题归结到\hyperref[fga3.iii-5-theorem-5.1]{定理~\ref*{fga3.iii-5-theorem-5.1}}, 其证明类似于\emph{Séminaire Chevalley}中代数商群的构造.
\end{proof}
\end{theorem}
\par{}归纳如下:
\begin{scholium}{}\label{fga3.iii-6-scholium}\skipsinglepar
\par{}设\(X\)在\(S\)上拟射影, 且\(S\)局部Noether. 则给出从\(X\)到某个\(S\)-预概形\(Y\)的一个固有、忠实平坦且满的态射\(f\colon  X\to  Y\), 等价于给出\(X\)上的一个等价关系\(R\), 使得\(p_1\colon  R\to  X\)固有且平坦.
\end{scholium}
\par{}同一方法还给出以下结果:
\begin{theorem}{6.2}\setnumber{6.2}\label{fga3.iii-6-theorem-6.2}\skipsinglepar
\par{}设\(S\)是Noether预概形, \(X\)是在\(S\)上有限型的预概形, \(\mathcal {R}\)是\(S\)-预概形\(X\)上的预等价关系. 假设
\begin{enumerate}
      
 \item[a.]\(p_1\colon  R_1\to  X\)平坦且有限型;

      
 \item[b.]态射\(R_1\to  X\times _S X\)拟有限, 即其纤维有限.

    
\end{enumerate}. 则\(X\)中存在一个关于\(\mathcal {R}\)\emph{饱和}的\emph{稠密}开子集\(U\), 使得: 
\begin{enumerate}
      
 \item[i.]若\(\mathcal {R}_U\)是\(\mathcal {R}\)在\(U\)上诱导的预等价关系, 则\(U/\mathcal {R}_U\)存在, 并在\(S\)上有限型.

      
 \item[ii.]典范态射\(U\to  U/\mathcal {R}_U\)满且开, 其纤维是\(\mathcal {R}_U\)的集合论等价类 (因此\(U/\mathcal {R}_U\)是\(U\)关于\(\mathcal {R}_U\)所定义的集合论等价关系的拓扑商空间). 最后, 态射\((\mathcal {R}_U)_1\to  U\times _{U/\mathcal {R}_U}U\)是满的.

      
 \item[iii.]若\(\mathcal {R}\)来自等价关系, 则可以假设\(U\to  U/\mathcal {R}_U\)忠实平坦, 且\(\mathcal {R}_U\)有效.

    
\end{enumerate}
\end{theorem}
\par{}这是一个本质上具有"双有理"性质的结果.
\begin{remarks}{6.3}\setnumber{6.3}\label{fga3.iii-6-remarks-6.3}\skipsinglepar
\begin{enumerate}\item{}我不知道在\hyperref[fga3.iii-6-theorem-6.1]{定理~\ref*{fga3.iii-6-theorem-6.1}}和\hyperref[fga3.iii-6-theorem-6.2]{定理~\ref*{fga3.iii-6-theorem-6.2}}中, 假设(b)是否多余. 在实际的群的取商中, 这个假设迫使我们限于所有稳定子都是有限群的情形.

    \item{}可以问, 在没有任何平坦性假设时, 是否也有类似于\hyperref[fga3.iii-6-theorem-6.1]{定理~\ref*{fga3.iii-6-theorem-6.1}}和\hyperref[fga3.iii-6-theorem-6.2]{定理~\ref*{fga3.iii-6-theorem-6.2}}的结果. 我在这一方向还没有反例. 但是, 即使保留平坦性假设, 限于\(p_1\colon  R\to  X\)平坦拟有限 (但不有限) 且\(X\)仿射的等价关系, \(R\)仍可能无效: 只需取覆盖Nagata簇的仿射开子集上所诱导的等价关系 (在该簇上, 一个二元群以"非容许"的方式作用).
\end{enumerate}
\end{remarks}
\chapter{应用}\thispagestyle{fancy}\label{fga3.iii-7}
\par{}如引言所述, \hyperref[fga3.iii-6-theorem-6.1]{定理~\ref*{fga3.iii-6-theorem-6.1}}最重要的应用是构造Picard概形以及解决其他各种"模"问题; 我们以后还会回到这些问题.
\par{}我们得到Shimura的下述结果的一个简单证明:
\begin{proposition}{7.1}\setnumber{7.1}\label{fga3.iii-7-proposition-7.1}\skipsinglepar
\par{}设\(A\)是在分式域为\(K\)的离散赋值环\(V\)上定义的阿贝尔概形. 则每个与\(A\otimes _V K\)的某个商同源的\(K\)上阿贝尔概形\(B'\), 都"在\(V\)处有好约化", 即它同构于某个\(B\otimes _V A\), 其中\(B\)是\(V\)上的阿贝尔概形 (回顾一下, 它本质上唯一).
\begin{proof}
\par{}可以假设\(B'\)是\(A_K\)关于一个子群概形\(C'\)的商. (注意, \(C'\)一般不"约化", 即其局部环可能有幂零元.) 考虑\(A\)中由\(C'\)的"闭包"给出的闭子概形\(C\), 即使\(C_K\)包含\(C'\)的\(A\)的最小闭子概形. 则\(C_K=C'\); 又因\(V\)是离散赋值环, 容易推出\(C\)是\(A\)在\(V\)上的子群概形. 由于\(A\)在\(\operatorname {Spec}(V)=S\)上固有, \(C\)也固有. 此外, \(A\)在\(S\)上射影. 于是可以应用\hyperref[fga3.iii-6-theorem-6.1]{定理~\ref*{fga3.iii-6-theorem-6.1}}构造\(A/C=B\), 这就是所求的\(B\).
\end{proof}
\end{proposition}
\par{}最后, 用基本上已知的论证可从\hyperref[fga3.iii-6-theorem-6.2]{定理~\ref*{fga3.iii-6-theorem-6.2}}得出以下结果:
\begin{theorem}{7.2}\setnumber{7.2}\label{fga3.iii-7-theorem-7.2}\skipsinglepar
\par{}设\(S\)是Artin环的谱, \(F\)和\(G\)是在\(S\)上有限型的群概形, \(u\colon  F\to  G\)是\(S\)上的群概形同态. 假设
\begin{enumerate}
      
 \item[i.]\(F\)在\(S\)上平坦;

      
 \item[ii.]\(u\)的核有限.

    
\end{enumerate}. 在这些条件下, 商概形\(G/F\)存在; 典范态射\(G\to  G/F\)满且开, 其纤维是\(F\)在\(G\)上的右作用定义的集合论等价类. 最后, 若\(u\)是单态射, 则态射\(G\to  G/F\)平坦, 而态射\(G\times  F\to  G\times _{(G/F)}G\)是同构; 换言之, \(G\)是\(G/F\)上以\(F\) (右作用) 为结构群的主齐性空间, 更准确地说, 结构群是视为\(G/F\)上群概形的\(F\times _S(G/F)\) (参见\hyperref[fga3.i-b.6]{[1], B, §6}).
\end{theorem}
\begin{corollary}{7.3}\setnumber{7.3}\label{fga3.iii-7-corollary-7.3}\skipsinglepar
\par{}在这些条件下, \(G\)在\(S\)上平坦的充要条件是\(G/F\)在\(S\)上平坦. 若此条件满足, 则关于\(F\)取商与基底\(S\)的任意扩张交换; 若\(F\)是\(G\)的不变子群, 则\(G/F\)具有\(G\)关于\(F\)的\emph{商群}结构.
\end{corollary}
\par{}当\(S\)是域的谱时情形尤为简单, 因为此时每个\(S\)-预概形自动在\(S\)上平坦. 我们得到:
\begin{corollary}{7.4}\setnumber{7.4}\label{fga3.iii-7-corollary-7.4}\skipsinglepar
\par{}设\(F\)和\(G\)是域\(k\)上的有限型群概形, \(u\colon  F\to  G\)是\(k\)-群同态. 则\(u\)可分解为\(F\to  F'\to  G\), 其中\(F\to  F'\)是关于\(F\)的闭子群\(\operatorname {Ker} u\)取商所得的同态, 而\(F'\to  G\)是一个闭浸入的群同态. 商\(G/F=G/F'\)存在. 通常的形式运算 (例如Noether同构定理中的运算) 在\(k\)上的代数群之间成立.
\end{corollary}
\par{}这个结果使我们可以统一处理经典意义的代数群 (即在\(k\)上不可约且在\(k\)上光滑的群) 的取商, 以及Cartier所考虑的关于"无穷小"子群的取商. 将Cartier在Dieudonné关于形式群的工作之后引入的"超代数"视为\(k\)上形式概形范畴中的群, 并在适当情形 (当它们对应于\(k\)上秩有限的超代数时) 视为\(k\)上有限的代数群, 是有利的.
\chapter{一个猜想}\thispagestyle{fancy}\label{fga3.iii-8}
\begin{remark}{}\label{fga3.iii-8-remark-i}\skipsinglepar
\par{}\emph{[Comp.]} 现在看来, 本节的猜想即使对于特征\(0\)的域上的非奇异簇也不成立: 无论商的存在性还是其拟射影性都可能失败, 甚至在\(G\)的作用具有闭图时也是如此.
\end{remark}
\par{}这里的猜想源于一个需要: 要知道如何对作用在"Hilbert概形"的某些子概形上的射影群取商 (在概形理论中, 这些"Hilbert概形"用来替代Chow簇).
\par{}设\(S\)是预概形, \(n\)是整数. 对每个\(S\)上的预概形\(S'\), 配以元素取值于\(\mathcal {O}_{S'}\)的截面环的可逆\((n\times  n)\)矩阵群\(\operatorname {GL}(n,\Gamma (S',\mathcal {O}_{S'}))\). 这样得到一个关于\(S'\)的反变函子, 很容易证明它可表示; 因而它对应于\(S\)上的一个群概形 (并且在\(S\)上仿射), 记作\(\operatorname {GL}(n)_S\). 它的构造与基底变换相容, 所以实际上所有这些均来自\(\mathbb {Z}\)上的一个群概形, 记作\(\operatorname {GL}(n)\). 群\(\operatorname {GL}(1)\)称为\emph{乘法群}, 常记作\(\operatorname {G_m}\), 它对应于函子\(S\mapsto \Gamma (S,\mathcal {O}_S)^*\), 后者是\(S\)上的"单位"群. 有一个显然的同态\(\operatorname {GL}(1)\to \operatorname {GL}(n)\), 且很容易构造其商群, 记作\(\operatorname {GP}(n-1)\), 称为\emph{\(\mathbb {Z}\)上次数为\((n-1)\)的射影群}. 它表示把\(S\)对应于层\(\operatorname {\mathscr {G}\kern  -2pt\mathscr {L}}(n)_S/\operatorname {\mathscr {G}\kern  -2pt\mathscr {L}}(1)_S\)的截面群的函子; 这里\(\operatorname {\mathscr {G}\kern  -2pt\mathscr {L}}(n)_S\)表示\(\operatorname {GL}(n)_S\)在\(S\)上的截面芽层. (注意, \(\operatorname {GP}(n-1)_S\)在\(S\)上的截面一般不来自\(\operatorname {GL}(n)_S\)在\(S\)上的截面!) 还可证明, 至少当\(S\)为Noether时, \(\operatorname {GP}(n-1)\)也表示函子\(S\mapsto \operatorname {Aut}_S(\mathbb {P}_S^{n-1})\) (其中\(\mathbb {P}_S^{n-1}\)是在\(S\)上相对维数为\((n-1)\)的射影型概形). 它正是这样出现在模理论中的.
\par{}设\(S\)是一个Noether概形, 也可以假设它仿射; 设\(X\)是一个拟射影的\(S\)-预概形, 带有关于\(S\)非常丰沛的可逆层\(\mathscr {L}\). 假设群\(G=\operatorname {GP}(n)_S\)同时作用于\(X\)及\(\mathscr {L}\) (其作用与\(X\)上的作用相容), 并在\(S\)上\emph{自由地}作用.
\begin{conjecture}{8.1}\setnumber{8.1}\label{fga3.iii-8-conjecture-8.1}\skipsinglepar
\par{}在上述条件下: 
\begin{enumerate}\item{}\(G\)定义的等价关系有效, 商\(Y=X/G\)在\(S\)上有限型, 且典范态射\(f\colon  X\to  Y\)平坦且满 (因而\(X\)成为\(Y\)上的齐性主丛, 群为\(G\times _S Y=\operatorname {GP}(n)_Y\)).

      \item{}设\(\mathscr {L}'\)是通过\(f\)下的"忠实平坦下降"从\(\mathscr {L}\)在\(Y\)上诱导的可逆层 (参见\hyperref[fga3.i-b.1-theorem-1]{[1], B, 定理1}). 则\(\mathscr {L}'\)在\(Y\)上关于\(S\)"预丰沛", 即存在一个整数\(m\)和从\(Y\)到某个适当射影型概形\(\mathbb {P}_S^N\)的拟有限态射, 使\((\mathscr {L}')^{\otimes  m}\)同构于\(\mathcal {O}_{\mathbb {P}_S^N}(1)\)的逆像.
\end{enumerate}
\end{conjecture}
\par{}应注意, 即使\(X\)在\(S\)上分离, \(Y\)也可能不在\(S\)上分离 (这种情形出现在完全不病态的"模问题"中). 如果(1)成立, 则\(Y\)分离当且仅当\(G\)所定义的等价关系具有闭图, 即\(G\times  X\to  X\times  X\)的像闭 (从而它是闭浸入). 若\(Y\)分离, 则\(\mathscr {L}'\)在\(Y\)上关于\(S\)预丰沛当且仅当它丰沛, 即某个适当张量幂定义一个射影浸入. 在引言提到的模问题中, 可以证明最终得到的等价关系确实具有闭图.
\begin{remarks}{8.1}\setnumber{8.1}\label{fga3.iii-8-remarks-8.2}\skipsinglepar
\par{}我们假设\(G=\operatorname {GP}(n)_S\), 是为了给出具体例子, 也因为这在目前实际应用中是最重要的情形. 对\(G\)作的合理假设, 似乎应是\(G\)为某个Tohokû群在\(S\)上的一种"形式" (Chevalley也已给出这些群在整数环上的构造). 关于上述猜想, 我所知道的唯一肯定事实如下: \emph{设\(X\)是特征\(0\)的域\(k\)上的仿射概形, 群\(\operatorname {GL}(n)_k\)或\(\operatorname {GP}(n-1)_k\)在其上自由作用. 则\(G\)定义的等价关系有效, 商\(X/G\)仿射, 且态射\(X\to  X/G\)平坦且满.} 证明使用下面的事实 (目前仅在特征\(0\)时得到证明): 让\(G\)作用于其仿射环\(A\) (此环为\(G\)的仿射环, 视为\(k\)上的向量空间), 则\(G\)的平凡表示仅出现一次 (在\(G\)下稳定的某个\(k\)上有限维向量子空间的合成列中). 系统地运用\(G\)的线性表示理论, 似乎可能证明这个猜想, 至少在基底是域时如此. 基底不再是域时, 作者一无所知.
\end{remarks}
\begin{remark}{}\label{fga3.iii-8-remark-ii}\skipsinglepar
\par{}\emph{[Comp.]} 正如下一讲末尾所指出的, 上述猜想确定是错误的. 上一注中的"肯定事实"似乎由不同作者同时证明 (Nagata、Rosenlicht、Grothendieck等).
\end{remark}

\section*{参考文献}

\begin{enumerate}
\item[{[1]}]\setnumber{1}\label{fga3.iii-bib-1} GROTHENDIECK (Alexander). --- Techniques de descente et théorèmes d'existence en géométrie algébrique, I: Généralités, \emph{Séminaire Bourbaki}, 1959/60, n\textsuperscript{o} 190, 29页.
\item[{[2]}]\setnumber{2}\label{fga3.iii-bib-2} GROTHENDIECK (Alexander). --- Techniques de descente et théorèmes d'existence en géométrie algébrique, II: Le théorème d'existence en théorie formelle des modules, \emph{Séminaire Bourbaki}, 1959/60, n\textsuperscript{o} 195, 22页.
\item[{[3]}]\setnumber{3}\label{fga3.iii-bib-3} GROTHENDIECK (A.) 与 DIEUDONNÉ (J.). --- \emph{Éléments de géométrie algébrique}, I: Le langage des schémas. Paris, Presses universitaires de France, 1960 (Institut des hautes Études scientifiques, Publications mathématiques, 4).

\end{enumerate}
